% TDSC core protocol text: System Model, Protocol, and Security Analysis.

%-------------------------------------------------------------------------------
\section{System Model and Preliminaries}
\label{sec:model}
%-------------------------------------------------------------------------------

This section defines the system boundary, authenticated receipt-order evidence,
fairness objective, and graph semantics used throughout AUTIG.  The
authenticated evidence state is the source of protocol truth; the leader's
UTIG is a derived acceleration structure.  Section~\ref{sec:protocol}
specifies the deterministic state transitions, extraction policy, and
verification rules.

\subsection{System Architecture}
\label{sec:system-architecture}

We consider a system composed of a static set of $n$ replicas,
$\mathcal R=\{r_1,\ldots,r_n\}$, which together provide a fair ordering
service.  At most $f$ replicas are Byzantine.  A replica is \emph{correct} if
it remains uncorrupted throughout the execution, and
$\mathcal R_C\subseteq\mathcal R$ denotes the set of correct replicas.  We
consider static membership within one protocol execution.  The hosting BFT
satisfies its standard resilience condition; we assume $n\ge3f+1$.  AUTIG
additionally requires Eq.~\eqref{eq:integer-feasibility} for
batch-order-fairness.

AUTIG is an ordering layer rather than a replacement for Byzantine consensus.
One replica acts as the \emph{order leader}: it collects authenticated
receipt-order evidence, materializes the dependency graph, and proposes a
fair-order fragment.  The other replicas act as followers that validate the
proposal before voting in the hosting BFT protocol.  The BFT layer totally
orders valid AUTIG records---fragments, checkpoints, and handoffs---or their
data-available digests, and provides agreement and finality.  A locally
verified fragment is therefore only a candidate.  Its evidence state and
output order become authoritative only after BFT commitment.

Fragments are indexed by sequence numbers $r=1,2,\ldots$; below, round $r$
refers to processing fragment sequence $r$.  The hosting BFT maintains the
order-leader authorization context
$\mathsf{AuthCtx}=\langle\mathsf{epoch},\leaderpk\rangle$.  A committed handoff
updates this context without changing the latest committed AUTIG evidence
state.  For fragment $r$, $(\mathsf{epoch}_r,\leaderpk)$ are the values in the
context authorizing that proposal.  A replacement leader resumes from the
latest committed evidence state and authorization context.

\subsection{Network and Adversarial Model}
\label{sec:network-adversary}

The network is partially synchronous~\cite{dwork1988consensus}.  Before an
unknown global stabilization time (GST), messages may be delayed arbitrarily.
After GST, messages between correct replicas are delivered within a bounded
delay $\Delta$.  The delay bound is required only for progress. 
We assume that the hosting BFT eventually installs a correct order leader
and continues to commit valid AUTIG fragments derived from the latest committed
state.

The adversary is computationally bounded and may adaptively corrupt up to $f$
replicas in total.  Corruptions are non-mobile: once corrupted, a replica
remains Byzantine.  Byzantine replicas may equivocate, omit or delay messages,
report arbitrary receipt orders, and collude with Byzantine clients.  The
adversary may also schedule network delivery subject to partial synchrony after
GST and can therefore influence the first-receipt orders observed by correct
replicas.

Clients may submit duplicate, conflicting, or strategically timed
transactions.  A transaction is \emph{valid} if it satisfies the deterministic
admission predicate of the replicated application.  The fairness guarantee
applies only to valid transactions whose complete contents become available
to all correct replicas.

\subsection{Cryptographic and Data-Availability Assumptions}
\label{sec:cryptographic-assumptions}

Every replica has a certified signing key, and signatures are existentially
unforgeable under chosen-message attack.  The hash function $H$ is collision
resistant, and $\Enc(\cdot)$ is a canonical encoding for every protocol object.
For protocol objects $(X_1,\ldots,X_k)$, define
$\Commit(X_1,\ldots,X_k;\mathsf{dom})
=H(\mathsf{dom}\parallel\Enc(X_1,\ldots,X_k))$, where $\mathsf{dom}$ is a
domain-separation tag.

Transaction identifiers are content-bound.  For canonical transaction bytes
$\mathsf{tx}_u$, AUTIG defines
$u=\Commit(\mathsf{tx}_u;\mathsf{AUTIG/transaction})$.  An application-native
identifier, when present, is encoded inside $\mathsf{tx}_u$ and is therefore
bound by $u$.  While $u$ remains live, replicas retain the consensus-level
content address $\mathsf{TxRef}_r[u]=u$; the referenced contents remain
retrievable under the data-availability assumption below.  Storage locators
are local and do not enter consensus state.

AUTIG relies on the hosting BFT's external-validity and data-availability
interface.  A correct replica votes for a fragment only after obtaining the
complete fragment and the transaction contents referenced by newly admitted
evidence, and after checking the application's validity predicate.  If the BFT
orders a digest rather than the bytes themselves, its availability rule must
ensure that the corresponding fragment remains retrievable.  Checkpoints apply
an analogous availability rule to the complete live-evidence recovery package
as specified in Section~\ref{sec:checkpoint-recovery}.

\subsection{Authenticated Receipt-Order Evidence}
\label{sec:receipt-order-evidence}

The fairness premise concerns actual first-receipt order at correct replicas.
For distinct valid transactions $u$ and $v$, we write $u\prec_i v$ if replica
$i$ receives the complete contents of $u$ before those of $v$.  Replica
$i$ validates received transactions and records their identifiers in an append-only observation log
$O_i[1],O_i[2],\ldots$ according to their first-receipt order. A correct
replica records each valid identifier only at its first local position and
reports signed contiguous extensions of this log. 
Therefore, a committed prefix reported by a correct replica and containing
$v$ also contains every valid transaction that the same replica received before
$v$.  Section~\ref{sec:prefix-consistent-localorders} gives the extension
format and admissibility checks.

Let $\mathsf{Pos}_r(i,u)\in\mathbb N\cup\{\bot\}$ be the unique first
position of live transaction $u$ in the authenticated prefix of replica $i$'s
log represented by evidence-state version $r$.  Its cumulative visibility
count is $c_r(u)=|\{i:\mathsf{Pos}_r(i,u)\ne\bot\}|$.
For an ordered transaction pair, the cumulative precedence weight is
\begin{equation}
 W_r(u,v)=\bigl|\{i:\mathsf{Pos}_r(i,u)\ne\bot,\,
 \mathsf{Pos}_r(i,v)\ne\bot,\,
 \mathsf{Pos}_r(i,u)<\mathsf{Pos}_r(i,v)\}\bigr|.
 \label{eq:position-weight}
\end{equation}
These values count distinct replica logs, not LocalOrder messages.  Because a
committed log has one first position per identifier, replica $i$ contributes at
most once to an unordered pair and to at most one of $W_r(u,v)$ and
$W_r(v,u)$.  Retransmission therefore cannot amplify evidence, and the
per-replica positions uniquely determine every cumulative weight.

Hatted symbols, such as $\widehat W_r$ and $\widehat A_r$, denote values
derived from the candidate pre-finalization state of fragment $r$.  Unhatted
symbols denote committed values unless the state version is fixed explicitly
by context.

The set $\mathsf{Live}_r$ contains admitted transactions that remain
unfinalized.  Finalized identifiers remain queryable through
$\mathsf{Done}_r$, which records their committed fragment and fairness-batch
index.  A finalized identifier may appear in a later contiguous log extension
for continuity, but it is not reactivated and does not contribute new
live-pair evidence.

\subsection{Fairness and Protocol Objectives}
\label{sec:fairness-objective}

AUTIG combines the agreement supplied by the hosting BFT with three ordering-
layer objectives.  First, \emph{evidence consistency} requires correct replicas
that commit the same fragment sequence to derive the same authenticated
receipt-order state.  Second, \emph{safe-fragment validity} requires an
accepted fragment to contain only valid, data-available transactions and to
finalize exactly the deterministic safe closure of the authenticated
pre-finalization state.  Third, \emph{batch-order-fairness} constrains the
relative fairness-batch indices assigned to transactions.  Section~\ref{sec:analysis}
establishes these properties separately.

Following the honest-replica parameterization of
Themis~\cite{kelkar2023themis}, AUTIG sets the required support to a
$\gamma$-fraction of the guaranteed correct-replica population $n-f$, where
$\gamma\in(1/2,1]$, and uses $q_h=\lceil\gamma(n-f)\rceil$.
Let the committed output be partitioned into consecutive fairness batches, and
write $b(u)$ for the index assigned to transaction $u$.  AUTIG satisfies
$\gamma$-batch-order-fairness if, for every pair of distinct valid transactions
$x$ and $y$ that are eventually received by all correct replicas, whenever
$y$ is assigned a fairness-batch index,
\begin{equation}
 \left|\{i\in\mathcal R_C:x\prec_i y\}\right|\ge q_h
 \quad\Longrightarrow\quad
 x\text{ is assigned and }b(x)\le b(y).
 \label{eq:batch-order-fairness}
\end{equation}
In particular, finalizing $y$ under the receive-order antecedent requires $x$
to be finalized in the same or an earlier fairness batch.  The property does
not impose a strict order on transactions assigned to the same batch.

AUTIG uses the following integer thresholds:
\begin{equation}
 T_{\mathrm{solid}}=n-2f,
 \qquad
 T_{\mathrm{non\mbox{-}blank}}=T_{\mathrm{edge}}
   =n-q_h+1.
 \label{eq:integer-thresholds}
\end{equation}
The fairness proof requires
\begin{equation}
 2q_h\ge n+2f+1.
 \label{eq:integer-feasibility}
\end{equation}
When $\gamma(n-f)$ is integral, Eq.~\eqref{eq:integer-thresholds} gives
$T_{\mathrm{edge}}=n(1-\gamma)+\gamma f+1$, and
Eq.~\eqref{eq:integer-feasibility} is equivalent to
$n>2f(\gamma+1)/(2\gamma-1)$.
We use Eqs.~\eqref{eq:integer-thresholds} and
\eqref{eq:integer-feasibility} when $\gamma(n-f)$ is non-integral.
Together with $n\ge3f+1$, Eq.~\eqref{eq:integer-feasibility} implies
$T_{\mathrm{non\mbox{-}blank}}\le T_{\mathrm{solid}}$, so the three state
intervals below are well defined.

Define $\textsc{DetermineState}(c)$ to return $\Solid$ when
$c\ge T_{\mathrm{solid}}$, $\Shaded$ when
$T_{\mathrm{non\mbox{-}blank}}\le c<T_{\mathrm{solid}}$, and $\Blank$
otherwise.  We write $S_r(u)=\textsc{DetermineState}(c_r(u))$.  The active set
induced by cumulative evidence is
$A_r=\{u\in\mathsf{Live}_r:S_r(u)\ne\Blank\}$.
Because $c_r(u)$ is cumulative while $u$ remains live, its state can only
progress from $\Blank$ to $\Shaded$ to $\Solid$.  Extraction and verification
operate on $A_r$.

\subsection{Graph-Based Ordering Preliminaries}
\label{sec:graph-preliminaries}

AUTIG represents cumulative precedence evidence as a directed dependency
graph.  Let $u\triangleleft v$ denote a public total order over transaction
identifiers, obtained by lexicographically comparing their canonical encodings.
For distinct transactions, $\mathsf{EdgePred}_r(u\to v)$ holds if and only if
$u,v\in A_r$, $W_r(u,v)\ge T_{\mathrm{edge}}$, and either
$W_r(u,v)>W_r(v,u)$ or both $W_r(u,v)=W_r(v,u)$ and $u\triangleleft v$.
The weight threshold excludes insufficiently observed relations, the comparison
retains the stronger supported direction, and the public identifier order
resolves an exact tie without leader discretion.  The active-set conjunct
excludes Blank transactions, and every weight contribution is derived from a
replica log containing both endpoints.

The protocol-defined dependency graph at sequence $r$ is
$G_r^\star=(A_r,E_r^\star)$, where
$(u,v)\in E_r^\star$ if and only if
$\mathsf{EdgePred}_r(u\to v)$.  The order leader may maintain an incremental
materialization of this graph as $G_{\mathrm{utig}}$.  The authenticated
positions and deterministic edge predicate uniquely define $G_r^\star$;
followers may evaluate individual predicates from local cumulative evidence
without storing $E_r^\star$.
For candidate state $\widehat\Sigma_r$, let
$\widehat{\mathsf{EdgePred}}_r$ denote the same rule evaluated with
$\widehat A_r$ and $\widehat W_r$; it uniquely defines
$\widehat G_r^\star$.

For a directed graph $G=(V,E)$, $\mathsf{Anc}_G(u)$ contains $u$ and every
vertex with a directed path to $u$.  A strongly connected component (SCC) is a
maximal set of mutually reachable vertices; contracting every SCC yields the
acyclic condensation graph.  A set $X\subseteq V$ is \emph{down-closed} if
$u\in X$ implies $\mathsf{Anc}_G(u)\subseteq X$.  Down-closure is the relevant
finalization condition: if a finalized vertex had an unfinalized incoming
ancestor, the committed output could place a supported predecessor too late.

Each finalized SCC receives one global fairness-batch index.  Transactions in
the same SCC therefore satisfy Eq.~\eqref{eq:batch-order-fairness} by sharing a
batch, while the condensation order constrains different SCCs.  AUTIG uses a
public deterministic serialization within an SCC only to supply the execution
interface; this serialization does not assert a stronger intra-batch fairness
property.  Section~\ref{sec:safe-extraction} defines which down-closed SCCs may
be finalized, and Section~\ref{sec:proof-carrying-verification} defines how a
graph-free follower verifies that choice.

\paragraph{Progress scope.}
Evidence consistency and safe-fragment validity hold for every BFT-committed
history under the stated system and cryptographic assumptions.
Batch-order-fairness additionally requires
Eq.~\eqref{eq:integer-feasibility}.  Progress assumes eventual dissemination
of valid transactions and eventual admission of every pending position in a
correct replica's log.  Together with the correct-leader and BFT-liveness
assumptions above, these conditions imply eventual finalization for any valid
transaction whose live ancestor set eventually becomes finite and stops
expanding.  Section~\ref{sec:discussion-scope} discusses this condition.

%-------------------------------------------------------------------------------
\section{The AUTIG Protocol}
\label{sec:protocol}
%-------------------------------------------------------------------------------

%-------------------------------------------------------------------------------
\subsection{Protocol Overview}
\label{sec:protocol-overview}
%-------------------------------------------------------------------------------

\paragraph{State and roles.}
AUTIG separates authenticated evidence replication from graph materialization.
Every replica applies committed, prefix-consistent receipt logs and maintains
the resulting live positions.  The order leader additionally materializes the
derived dependency graph as a persistent UTIG and incrementally updates its
affected region.  Followers remain graph-free: they derive pair predicates
from authenticated positions and verify a structural certificate without
maintaining the UTIG edge set, SCC decomposition, or condensation DAG.

\paragraph{Round workflow.}
Each protocol round proceeds as follows.  Replicas extend their signed
observation logs; the order leader applies an admitted batch, refreshes the
UTIG, extracts a safe down-closed set, and constructs its certificate.
Followers replay the same evidence transition and verify the certificate
against locally derived weights.  This produces a candidate post-state, which
becomes authoritative only when the hosting BFT commits the fragment.
Checkpoints preserve the complete live-evidence state needed for recovery and
order-leader handoff.

\paragraph{Safe extraction.}
A transaction is finalizable only if it is Solid and all its graph ancestors
are Solid.  A blocker forest explains each excluded Solid transaction, and a
complete incoming-frontier check proves that the proposed set is down-closed.
Thus followers verify the selected subgraph and its boundary without
materializing or extracting the complete active dependency graph.

%-------------------------------------------------------------------------------
\subsection{Prefix-Consistent LocalOrders}
\label{sec:prefix-consistent-localorders}
%-------------------------------------------------------------------------------

As defined in Section~\ref{sec:receipt-order-evidence}, replica $i$ records
valid transactions in the order in which their complete contents are first
received.  Signed contiguous extensions of this observation log make prefix
completeness enforceable during evidence replay.

For sequence $r$, $L_i^r$ contains
$(\mathsf{epoch}_r,r,i,\ell_i^{-},\ell_i^{+},\eta_i^{-},
\mathsf{chunk}_i^r,\eta_i^{+})$ and a signature $\sigma_i$ over these fields.
The chunk covers positions $\ell_i^{-}+1,\ldots,\ell_i^{+}$, and
$\eta_i^{+}=\Commit(\eta_i^{-},\ell_i^{-},\ell_i^{+},
\mathsf{chunk}_i^r;\mathsf{AUTIG/local\mbox{-}log})$.
Chunks are bounded by $lo\text{-}size$ but cannot skip positions.

Let $\mathsf{Seq}_{r-1}[i]$ and $\mathsf{Head}_{r-1}[i]$ be the position and
log head installed for $i$ by the preceding committed fragment.
A LocalOrder is admissible only if its signature, epoch, sender, and round are
valid,
$\ell_i^{-}=\mathsf{Seq}_{r-1}[i]$,
$\eta_i^{-}=\mathsf{Head}_{r-1}[i]$, its positions are contiguous, its
identifiers are distinct within the extension, and no listed identifier
outside $\mathsf{Done}_{r-1}$ already has a live position for $i$.  A batch
$\mathcal L_r$ contains exactly $n-f$ admissible extensions from distinct
senders.

Log validation is distinct from contribution to live evidence.  A delayed,
chain-valid extension may contain identifiers in $\mathsf{Done}_{r-1}$; such
entries preserve log continuity but neither reactivate transactions nor add
incident pair evidence.  Before assigning any other first position, a follower
obtains the digest-bound contents, verifies application validity, and stores
$\mathsf{TxRef}_r[u]$; missing or inconsistent contents reject the candidate.
Correct replicas retransmit pending chunks.  After a committed handoff that
preserves sequence $r$, they re-sign the unchanged next contiguous chunk under
the successor epoch; the committed $\mathsf{Seq}$ and $\mathsf{Head}$ remain
unchanged.  A correct leader uses a public round-robin sender schedule, and
insufficient input initiates order-leader handoff.  Under
Section~\ref{sec:graph-preliminaries}'s progress assumptions,
every finite correct-replica log position is eventually admitted.  The complete
admissibility and retransmission rules appear in
Appendix~\ref{app:expanded-protocol}.

%-------------------------------------------------------------------------------
\subsection{Authenticated Graph-Free Follower State}
\label{sec:authenticated-evidence-state}
%-------------------------------------------------------------------------------

A follower persists the authoritative positions $\mathsf{Pos}_r$, sequence
number $\mathsf{Seq}_r[i]$, and hash-chain head $\mathsf{Head}_r[i]$ for every
replica log.  Visibility, pair weights, transaction states, and the active set
are deterministic functions of these positions, as defined in
Sections~\ref{sec:receipt-order-evidence} and~\ref{sec:fairness-objective}.

The replicated state machine also maintains a queryable finalized-identifier
dictionary $\mathsf{Done}_r$, which maps each identifier to its committed round
and fairness-batch index, with authenticated root
$\mathsf{DoneRoot}_r=\Commit(\mathsf{Done}_r;
\mathsf{AUTIG/finalized\mbox{-}set})$.  Checkpoints use a canonical
$\mathsf{DoneRef}_r$: either an inline encoding of $\mathsf{Done}_r$ or an
application-checkpoint identifier paired with $\mathsf{DoneRoot}_r$.  Recovery
resolves either form to a queryable dictionary and verifies its root.

The follower state installed after committed fragment $r$ is
\begin{equation}
 \Sigma_r=\langle r,\mathsf{epoch}_r,\beta_r,\mathsf{Seq}_r,
 \mathsf{Head}_r,\mathsf{Live}_r,\mathsf{TxRef}_r,\mathsf{Pos}_r,
 \mathsf{Done}_r,\mathsf{DoneRoot}_r,h_r\rangle.
 \label{eq:follower-state}
\end{equation}
$\beta_r$ is the largest fairness-batch index assigned through fragment $r$
(with $\beta_0=0$).  The state field $\mathsf{epoch}_r$ records the context that
authorized fragment $r$; after a handoff, the current $\mathsf{AuthCtx}$ may
advance before the next evidence state is committed.  The logical view
$\mathsf{View}_r=\langle c_r,W_r,S_r,A_r\rangle
=\mathsf{Derive}(\Sigma_r)$ is uniquely determined by $\mathsf{Pos}_r$.
Implementations may cache this view, but it is non-authoritative.
$\mathsf{TxRef}_r$ binds each live identifier to its available contents.
Followers do not maintain the UTIG edge set, SCCs, condensation DAG, or
topological order.

The transition identifiers are
$\widehat h_r=\Commit(\mathsf{epoch}_r,r,h_{r-1},H(\Enc(\mathcal L_r));
\mathsf{AUTIG/prestate})$ and
$h_r=\Commit(\mathsf{epoch}_r,r,\widehat h_r,F_r,\mathsf{Part}_r;
\mathsf{AUTIG/poststate})$.  The $h$ field of $\widehat\Sigma_r$ is
$\widehat h_r$; the field of $\Sigma_r^{\mathsf{cand}}$ and committed
$\Sigma_r$ is $h_r$.  Checkpoint digests separately bind materialized recovery
state.  The hosting BFT supplies the committed-fragment digest
$D_r=H_{\mathcal F_r}$ (with a protocol-defined genesis value $D_0$) together
with the current $\mathsf{AuthCtx}$; neither context object is a field of
$\Sigma_r$.

%-------------------------------------------------------------------------------
\subsection{Deterministic Evidence Transition and UTIG Refresh}
\label{sec:evidence-transition}
%-------------------------------------------------------------------------------

Proposal verification and commitment follow the state transition:
\begin{equation}
 \Sigma_{r-1}\xrightarrow{\;\mathcal L_r\;}
 \widehat\Sigma_r
 \xrightarrow{\;(F_r,\mathsf{Part}_r)\;}
 \Sigma_r^{\mathsf{cand}}
 \xrightarrow{\;\textsf{BFT commit}\;}\Sigma_r.
 \label{eq:proposal-state-transition}
\end{equation}
$\widehat\Sigma_r$ is the candidate pre-finalization evidence state, and
$\Sigma_r^{\mathsf{cand}}$ is the candidate post-finalization state.  Both are
speculative; only $\Sigma_r^{\mathsf{cand}}$ is installed as $\Sigma_r$ after
the hosting BFT commits the fragment.

\textsc{ApplyEvidence} validates the $n-f$ extensions, installs each new
non-finalized first position exactly once, and derives the affected visibility
counts, pair weights, states, and active vertices.  It sets the candidate
version, epoch, and $h$ fields to $r$, $\mathsf{epoch}_r$, and
$\widehat h_r$, respectively.  Its normative per-position procedure appears in
Appendix~\ref{app:expanded-protocol}.

\begin{algorithm}[h!]
\caption{\textsc{BuildFragment}: Leader-Side Construction}
\label{alg:build-fragment}
\footnotesize
\Input{$\Sigma_{r-1}$, $\mathsf{View}_{r-1}$,
       $G_{\mathrm{utig},r-1}$, committed predecessor digest $D_{r-1}$,
       $\mathsf{AuthCtx}=\langle\mathsf{epoch}_r,\leaderpk\rangle$, $r$, and
       $\mathcal L_r$}
\Output{$\mathcal F_r$, $\Sigma_r^{\mathsf{cand}}$, and
        $\widehat G_{\mathrm{utig},r}$, or reject}
$(\widehat\Sigma_r,\widehat{\mathsf{View}}_r,
  \mathsf{TouchedNodes}_r,\mathsf{TouchedPairs}_r)
  \gets\textsc{ApplyEvidence}(\Sigma_{r-1},\mathsf{View}_{r-1},
  (\mathsf{epoch}_r,r),\mathcal L_r)$; reject on failure\;
$\widehat G_{\mathrm{utig},r}\gets G_{\mathrm{utig},r-1}$\;
Add every newly active vertex to $\widehat G_{\mathrm{utig},r}$ and
re-evaluate all active pairs in $\mathsf{TouchedPairs}_r$\;
Compute the safe SCCs, their deterministic ranks, the partition
$\mathsf{Part}_r$, and certificate $\mathcal P_r$\;
$F_r\gets$ the SCCs in rank order, with identifiers sorted within each SCC\;
$\Sigma_r^{\mathsf{cand}}\gets
  \textsc{Finalize}(\widehat\Sigma_r,F_r,\mathsf{Part}_r)$\;
Construct and sign $\mathcal F_r$ with $H_{\mathcal F,\prev}=D_{r-1}$ as specified in
Section~\ref{sec:proof-carrying-verification}; \KwRet
$(\mathcal F_r,\Sigma_r^{\mathsf{cand}},\widehat G_{\mathrm{utig},r})$\;
\end{algorithm}

Pairs are touched whenever a new unique position changes either directional
weight and whenever a vertex becomes active; the latter step also inserts that
vertex into the candidate UTIG.  Re-evaluation therefore captures threshold,
dominance, and tie changes and yields
$\widehat G_{\mathrm{utig},r}=\widehat G_r^\star$, where
$\widehat G_r^\star=(\widehat A_r,\widehat E_r^\star)$ is defined by the
candidate evidence state.  Appendix~\ref{app:incremental-materialization}
proves this equivalence.

\textsc{Finalize} assigns the certified SCC batch indices, advances $\beta_r$,
moves $F_r$ from $\mathsf{Live}$ to $\mathsf{Done}$, recomputes
$\mathsf{DoneRoot}_r$, and removes finalized positions and $\mathsf{TxRef}$
entries while preserving all remaining live positions.  It computes $h_r$
from the transition definition rather than copying the header value.  Removing
$\mathsf{TxRef}[u]$ affects only AUTIG's live-evidence state; the hosting BFT or
application retains committed contents until execution and its applicable
checkpoint policy.  Later log occurrences of finalized identifiers are
continuity-only.

%-------------------------------------------------------------------------------
\subsection{Safe Finalization and Fairness Batches}
\label{sec:safe-extraction}
%-------------------------------------------------------------------------------

The leader runs extraction on $\widehat G_{\mathrm{utig},r}$, which equals the
protocol-defined $\widehat G_r^\star$.  In this subsection,
$\mathsf{Solid}_r=\{u\in\widehat A_r:\widehat S_r(u)=\Solid\}$ abbreviates the
candidate pre-finalization Solid set.
For $u\in\widehat A_r$, $\mathsf{Anc}_r(u)$ contains $u$ and every vertex
that reaches $u$ in $\widehat G_r^\star$.  AUTIG finalizes the set
\begin{equation}
 K_r^{\mathsf{safe}}
   =\{u\in\mathsf{Solid}_r:
          \mathsf{Anc}_r(u)\subseteq\mathsf{Solid}_r\}.
 \label{eq:safe-closure}
\end{equation}
This is the unique maximal subset of Solid transactions that is down-closed
in the complete active dependency graph.  Equivalently, after condensing
$\widehat G_r^\star$, an SCC is safe if every transaction in that SCC and every
transaction in each ancestor SCC is Solid.  Safety can be computed in one
topological pass: a component is unsafe if it contains a non-Solid vertex or
has an unsafe predecessor.  A Solid transaction with a Shaded ancestor
remains live until that ancestor becomes Solid.

The SCC boundaries define the fairness batches.  Let $\beta_{r-1}$ be the
largest index assigned by earlier committed fragments.  If
$C_1,\ldots,C_k$ are the SCCs of
$\widehat G_r^\star[K_r^{\mathsf{safe}}]$ in topological order, each
$u\in C_j$ receives $b(u)=\beta_{r-1}+j$, and
$\beta_r=\beta_{r-1}+k$.
An empty safe set leaves $\beta_r=\beta_{r-1}$.  The fragment therefore carries
a sequence of fairness batches, rather than treating the entire fragment as
one batch.  For execution, AUTIG serializes the transactions within $C_j$ by
increasing transaction identifier; this serialization is deterministic but
does not change their common fairness-batch index.  Across SCCs, Kahn's
algorithm repeatedly selects the zero-indegree component with the smallest
$\mathsf{cid}(C)=\min_{\triangleleft} C$.  The public rules determine
both the batch sequence and its execution serialization.

The structural certificate is
$\mathcal P_r=\langle\mathsf{Part}_r,\mathsf{InTree}_r,
\mathsf{OutTree}_r,\mathsf{Rank}_r,\mathsf{BlockForest}_r\rangle$.
$\mathsf{Part}_r$ maps $F_r$ to claimed SCCs.  Two trees rooted at each
component's smallest identifier certify strong connectivity, and
$\mathsf{Rank}_r$ records the deterministic condensation order.
$\mathsf{BlockForest}_r$ is a predecessor map rooted at Shaded vertices.  For
each non-root $v$, its stored parent $p(v)$ satisfies
$\widehat{\mathsf{EdgePred}}_r(p(v)\to v)$.  Following parent pointers from
every excluded Solid transaction reaches a Shaded root; equivalently, the
candidate graph contains a path from that root to the transaction.  Shared
paths keep the forest linear in its represented vertices.  Followers derive
all required pair weights from $\widehat{\mathsf{View}}_r$, so $\mathcal P_r$
carries no weight totals.  The witness forest need not be unique; the fragment
digest binds the leader's selected valid witness.
The leader computes the safe SCCs in one topological pass, applies the public
Kahn and identifier tie-breakers, and constructs the SCC trees and blocker
forest.  Expanded Algorithm~\ref{alg:safe-extract} is the normative
procedure.

Appendix~\ref{app:running-example} illustrates evidence derivation, safe
closure, structural certification, and follower verification in one example.

%-------------------------------------------------------------------------------
\subsection{Proof-Carrying Fragment Verification}
\label{sec:proof-carrying-verification}
%-------------------------------------------------------------------------------

The leader forms
$\hdr_r=\langle\mathsf{epoch}_r,r,\leaderpk,H_{\mathcal F,\prev},
h_{r-1},\widehat h_r,h_r,H_{\mathcal F_r}\rangle$ and
$\mathcal F_r=\langle\hdr_r,F_r,\mathcal L_r,\mathcal P_r,
\sigma_r^{\mathsf L}\rangle$, where
\begin{equation}
 H_{\mathcal F_r}=\Commit(\mathsf{epoch}_r,r,\leaderpk,
 H_{\mathcal F,\prev},h_{r-1},\widehat h_r,h_r,
 F_r,\mathcal L_r,\mathcal P_r;\domfrag).
 \label{eq:fragment}
\end{equation}
The order leader signs this digest as $\sigma_r^{\mathsf L}
=\mathsf{Sign}_{\mathsf{sk}_{\mathsf L}}(H_{\mathcal F_r})$ for the authorized
$\leaderpk$.  The digest binds the transition identifiers, evidence batch,
order, certificate, and predecessor fragment.  The BFT layer commits the full
fragment or its data-available digest.

\begin{algorithm}[h!]
\caption{\textsc{VerifyFragment}: Verify a Proof-Carrying Fragment}
\label{alg:verify-fragment}
\footnotesize
\Input{Committed state $\Sigma_{r-1}$; committed predecessor digest
       $D_{r-1}$; current $\mathsf{AuthCtx}$; candidate fragment $\mathcal F_r$}
\Output{$(accept,\Sigma_r^{\mathsf{cand}},h_r)$ or reject}
Require the header's $(\mathsf{epoch}_r,\leaderpk)$ to equal
$\mathsf{AuthCtx}$, require $r=\Sigma_{r-1}.r+1$, and verify
$\sigma_r^{\mathsf L}$ and
Eq.~\eqref{eq:fragment}; require
$h_{r-1}=\Sigma_{r-1}.h$ and $H_{\mathcal F,\prev}=D_{r-1}$\;
Derive $(\widehat\Sigma_r,\widehat{\mathsf{View}}_r)$ by
\textsc{ApplyEvidence}; reject on failure or a mismatching $\widehat h_r$\;
Require $F_r$ to contain distinct identifiers with
$\mathrm{set}(F_r)\subseteq\mathsf{Solid}_r$; using local $\widehat W_r$,
verify that $\mathsf{Part}_r$ is exactly the SCC partition of
$\widehat G_r^\star[\mathrm{set}(F_r)]$ and verify its trees, deterministic
condensation ranks, batch order, and serialization\;
Require $\neg\widehat{\mathsf{EdgePred}}_r(x\to y)$ for every
$(x,y)\in(\widehat A_r\setminus\mathrm{set}(F_r))
\times\mathrm{set}(F_r)$\;
Verify the blocker forest and require the parent pointers of every
excluded Solid transaction to reach a Shaded root\;
$\Sigma_r^{\mathsf{cand}}\gets
  \textsc{Finalize}(\widehat\Sigma_r,F_r,\mathsf{Part}_r)$\;
Verify its $h$ field equals the header's $h_r$\;
\KwRet $(accept,\Sigma_r^{\mathsf{cand}},h_r)$
\end{algorithm}

The tree and rank checks certify the SCC partition; the frontier scan proves
down-closure; and the blocker forest gives every excluded Solid transaction a
Shaded ancestor.  Together these checks establish
$\mathrm{set}(F_r)=K_r^{\mathsf{safe}}$, including the empty-set case, and the
deterministic rank and serialization rules fix its fairness batches.  The
normative field-level checks appear in
Algorithm~\ref{alg:verify-fragment-expanded}; Appendix~\ref{app:autig-security-details}
proves their soundness.

Verification does not mutate $\Sigma_{r-1}$.  Once the hosting BFT commits
$H_{\mathcal F_r}$, its context advances to $D_r=H_{\mathcal F_r}$; a follower
installs $\Sigma_r^{\mathsf{cand}}$ and releases $F_r$ for execution.  The
order leader likewise installs $\widehat G_{\mathrm{utig},r}$ with $F_r$ and
its incident edges removed as $G_{\mathrm{utig},r}$ only after that commit.

%-------------------------------------------------------------------------------
\subsection{Checkpoint and Recovery}
\label{sec:checkpoint-recovery}
%-------------------------------------------------------------------------------

Every $K_{\mathrm{cp}}$ committed fragments, replicas encode
$\mathsf{Snap}_r=\Enc(\mathsf{epoch}_r,r,\beta_r,\mathsf{Seq}_r,
\mathsf{Head}_r,\mathsf{Live}_r,\mathsf{TxRef}_r,\mathsf{Pos}_r,
\mathsf{DoneRef}_r,h_r)$ and set
$d_r=\Commit(\mathsf{Snap}_r;\mathsf{AUTIG/checkpoint\mbox{-}snapshot})$.  They commit
\begin{equation}
 \mathsf{CP}_r=\langle\mathsf{AuthCtx}_r,r,h_r,D_r,d_r,
                       \mathsf{AvailQC}_r\rangle,
 \label{eq:checkpoint}
\end{equation}
Here $\mathsf{AuthCtx}_r$ is the authorization context current when the
checkpoint is committed; after a handoff, it need not equal the epoch recorded
in $\Sigma_r$.  Moreover, $D_r=H_{\mathcal F_r}$ and $\mathsf{AvailQC}_r$ contains $n-f$
signatures on $(\mathsf{AuthCtx}_r,r,h_r,D_r,d_r)$.  A correct replica signs only its local
committed snapshot and retains it together with all data referenced by
$\mathsf{DoneRef}_r$ and $\mathsf{TxRef}_r$ until a successor checkpoint is
committed and available.  The certificate has
at least $n-2f$ correct holders, and the BFT availability layer retains the
complete committed AUTIG suffix and its referenced transaction contents after
the checkpoint.

A recovering replica validates a digest-matching snapshot, resolves
$\mathsf{DoneRef}_r$ and its authenticated root, validates all live-data
references, derives $\mathsf{View}_r$ from $\mathsf{Pos}_r$, and replays only
BFT-committed AUTIG suffix records.  Handoff records update
$\mathsf{AuthCtx}$; fragments advance the evidence state and predecessor digest
from $D_r$.  The snapshot preserves all live--live
positions, including pairs outside the last finalized prefix; recovery is
therefore exact.  A replacement leader materializes the UTIG from this state.
Algorithm~\ref{alg:recover-state} and Appendix~\ref{app:checkpoint-handoff}
give the complete procedure and proof.

%-------------------------------------------------------------------------------
\subsection{Pruning}
\label{sec:protocol-pruning}
%-------------------------------------------------------------------------------

Soft pruning may evict derived weights, edges, adjacency indexes, and
SCC/condensation metadata, provided that all authoritative live positions are
retained and every required live--live value can be reconstructed.  Hard
deletion applies only to evidence incident to a transaction after that
transaction appears in a BFT-committed $F_r$; $\mathsf{Done}_r$ prevents its
reactivation.  Evidence between two unfinalized transactions is never
hard-deleted.  Theorems~\ref{thm:pruning_safety} and~\ref{thm:hard-deletion},
with Appendix~\ref{app:pruning-retention}, establish that both operations
preserve future extraction and verification.

%-------------------------------------------------------------------------------
\subsection{Order-Leader Handoff}
\label{sec:order-leader-handoff}
%-------------------------------------------------------------------------------

Objective invalidity and silence use different evidence.  A fragment whose
bytes, LocalOrders, certificate, and referenced contents are available and that
bears a valid signature from the leader authorized by the current
$\mathsf{AuthCtx}$ but fails any deterministic verification check against
$(\Sigma_{r-1},D_{r-1})$ is independently verifiable $\mathsf{OrderFault}$
evidence and may be submitted to the hosting BFT for a replacement decision.
A fragment signed under a superseded authorization context is merely
inadmissible.  Fragment or transaction-content non-delivery is not publicly
verifiable and follows the timeout path.  After a
local timeout, $n-f$ matching signatures on
$\mathsf{OrderTimeout}(\mathsf{AuthCtx},r,h_{r-1},D_{r-1})$ form a handoff
certificate.  A BFT-committed handoff record names the successor
$\mathsf{AuthCtx}'=\langle\mathsf{epoch}+1,\leaderpk'\rangle$, with
$\leaderpk'$ selected by the hosting BFT's deterministic replacement rule, and atomically
replaces the current authorization context;
it changes neither $(\Sigma_{r-1},h_{r-1},D_{r-1})$ nor sequence $r$.  A late
fragment authorized by the replaced context is inadmissible.  The successor
resumes from the latest committed context, never from an uncommitted candidate.

%-------------------------------------------------------------------------------
\section{Security Analysis}
\label{sec:analysis}
%-------------------------------------------------------------------------------

AUTIG maintains five invariants across committed fragments: each replica log
extends one committed head; each replica contributes at most one direction to
each live pair; correct followers derive the same evidence state; uncommitted
candidate states cannot affect subsequent verification; and evidence between
two unfinalized transactions survives finalization, checkpointing, and
handoff.  We establish fairness in three steps.  First, prefix completeness
turns correct-replica receive-order observations into committed pair support.  Second,
the threshold condition makes that support both sufficient and strictly
dominant whenever the later transaction is Solid.  Third, safe-closure
verification places every required predecessor in the same or an earlier
fairness batch.

\begin{theorem}[Authenticated-state agreement]
\label{thm:state-agreement}
Correct replicas that commit the same fragment sequence derive identical
$\Sigma_r$, $h_r$, and derived views $\mathsf{View}_r$.
\end{theorem}

\begin{theorem}[Incremental-materialization equivalence]
\label{thm:incremental-equivalence}
If $G_{\mathrm{utig},r-1}=G_{r-1}^\star$, applying
Algorithm~\ref{alg:apply-evidence} and re-evaluating exactly
$\mathsf{TouchedPairs}_r$ yields
$\widehat G_{\mathrm{utig},r}=\widehat G_r^\star$.  After a committed
finalization removes $F_r$, the installed materialization satisfies
$G_{\mathrm{utig},r}=G_r^\star$.
\end{theorem}

\begin{theorem}[Safe-fragment validity]
\label{thm:safe-fragment}
If Algorithm~\ref{alg:verify-fragment} accepts $\mathcal F_r$, then
$F_r$ contains only application-valid, data-available transactions and
$\mathrm{set}(F_r)=K_r^{\mathsf{safe}}$.  Its partition is the SCC partition
of $\widehat G_r^\star[K_r^{\mathsf{safe}}]$, and its ranks and execution
order equal the public deterministic linearization.
\end{theorem}

\begin{theorem}[$\gamma$-Batch-Order-Fairness]
\label{thm:batch-fairness}
Assume Eq.~\eqref{eq:integer-feasibility}.  For distinct valid transactions
$x$ and $y$ that are eventually received by all correct replicas, if at least
$q_h$ correct replicas receive $x$ before $y$, then every committed AUTIG
execution that finalizes $y$ also finalizes $x$ and satisfies $b(x)\le b(y)$.
\end{theorem}

\begin{theorem}[Soft-pruning equivalence]
\label{thm:pruning_safety}
Let an unpruned leader and a soft-pruned leader start from the same committed
evidence state $\Sigma_r$.  The soft-pruned leader may evict any derived UTIG
cache, including edges, adjacency indexes, SCC metadata, condensation ranks,
topological order, and aggregate weights, but retains every authoritative
position and either retains or deterministically reconstructs every required
live--live pair value.  Rematerializing the evicted caches yields the same
$\widehat G_r^\star$, $K_r^{\mathsf{safe}}$, fairness-batch partition, and accepted
fragment as the unpruned execution.
\end{theorem}

\begin{theorem}[Hard-deletion safety]
\label{thm:hard-deletion}
After a committed fragment finalizes $u$, AUTIG may delete $u$'s live
positions and AUTIG live-data reference, and may evict all derived values
incident to $u$, without changing the fairness-batch assignment of any
unfinalized transaction or invalidating a future fairness obligation.
\end{theorem}

\begin{theorem}[Recovery equivalence]
\label{thm:recovery}
A replica that restores a digest-matching committed checkpoint and replays the
subsequent BFT-committed AUTIG records derives the same evidence state, latest
fragment digest, and order-leader authorization context as an uninterrupted
correct replica.
\end{theorem}

The proof intuition is as follows.  When $y$ is finalized, it is Solid.  If
the fairness antecedent holds and $x$ is still live, the quorum-intersection
bound establishes a strict edge $x\to y$.  Since the accepted prefix is
down-closed, $x$ belongs to the same prefix; SCC membership then assigns $x$
to either the same fairness batch as $y$ or an earlier one.  The same argument
shows that no late chunk can reveal a previously omitted fairness predecessor
of a finalized transaction, which permits finalized-incident deletion.
Checkpoint equivalence follows by induction over the deterministic evidence
transition.  Appendix~\ref{app:autig-security-details} gives the complete
proofs.

\begin{proposition}[Progress under stabilized ancestry]
\label{prop:stabilized-progress}
Assume eventual admission of every pending position in a correct replica's
log, eventual installation of a correct order leader after GST, and commitment
of its valid proposals by the hosting BFT.  Let $u$ be a valid live
transaction.  If its live ancestor set becomes finite and stops expanding,
and every transaction in that set is eventually disseminated to all correct
replicas, then $u$ is eventually included in a committed safe closure.
\end{proposition}

Every transaction in the resulting finite ancestor set eventually obtains at
least $n-f$ correct positions and therefore becomes Solid.  Once this occurs,
the entire ancestor set is safe, and verification of the maximal safe closure
requires its inclusion.  Without stabilized ancestry, new Shaded ancestors
may extend a dependency chain indefinitely.  Section~\ref{sec:discussion-scope}
states the resulting liveness scope.

%-------------------------------------------------------------------------------
\subsection{Complexity}
\label{sec:complexity}
%-------------------------------------------------------------------------------

Sparse positions require $O(n|\mathsf{Live}_r|)$ entries, and the aggregate
weight cache requires $O(|\mathsf{Live}_r|^2)$ counters in the worst case.  If
$N_i$ contains the new positions of replica $i$, evidence application touches
$O(|N_i||K_i^-|+|N_i|^2)$ pairs; no round-wide transaction-state scan is
required.  The SCC trees and ranks contain $O(|F_r|)$ records.  A blocker
forest contains $O(|V_{\mathsf{blk}}|)$ records, where
$|V_{\mathsf{blk}}|\le|\widehat A_r\setminus\mathrm{set}(F_r)|$; it is empty when every Solid
transaction is safe.  Verification reads
\[
 \frac{|F_r|(|F_r|-1)}{2}
 +( |\widehat A_r|-|F_r| )|F_r|+O(|V_{\mathsf{blk}}|)
\]
local pair values.  Thus, worst-case verification remains
$\Theta(|\widehat A_r|^2)$, although followers avoid persistent edge materialization,
SCC maintenance, and full active-graph extraction.

%-------------------------------------------------------------------------------
\subsection{Discussion and Scope}
\label{sec:discussion-scope}
%-------------------------------------------------------------------------------

AUTIG replicates authenticated receipt-order evidence at followers and does
not eliminate evidence-state maintenance.  Its structural certificate is
linear in the represented finalization and blocker subgraphs, while a verifier
may still inspect quadratically many local pair values in the worst case.  The
fairness definition assigns one batch index to every SCC; deterministic
serialization removes leader choice within that SCC but does not provide a
stronger intra-batch fairness guarantee.  The safe-closure rule provides
batch-order-fairness; under continuous arrivals, it provides progress once a
transaction's ancestry stabilizes.  Expiring an incoming edge solely because
its destination is old would define a different, bounded-fairness property: the
discarded edge may still be required by the receive-order antecedent.
Obtaining Themis-style standard liveness would require an additional
bounded-domain or deferred-ordering mechanism, such as committed ordering
frames and batch unspooling.  Designing and evaluating that composition is
outside the scope of this work.  The protocol analyzed here retains safe closure
and consequently establishes the stabilized-ancestry progress property of
Proposition~\ref{prop:stabilized-progress}.

%-------------------------------------------------------------------------------
\appendices
\section{Notation}
\label{app:notation}
%-------------------------------------------------------------------------------

Table~\ref{tab:notation} collects the notation used by the protocol and its
analysis.  Symbols carrying a hat refer to the candidate pre-finalization
state derived from the committed predecessor.  Only the corresponding
post-finalization candidate is installed after BFT commitment.

\begin{table*}[t]
\centering
\caption{Principal notation.}
\label{tab:notation}
\renewcommand{\arraystretch}{1.08}
\begin{tabular}{p{0.22\textwidth}p{0.72\textwidth}}
\hline
\textbf{Symbol} & \textbf{Meaning} \\
\hline
$n,f,\mathcal R,\mathcal R_C$ & Number of replicas, Byzantine bound, replica set, and correct-replica set. \\
$\gamma,q_h$ & Support fraction relative to the guaranteed correct-replica population $n-f$ and its integer threshold $q_h=\lceil\gamma(n-f)\rceil$. \\
$T_{\mathrm{solid}},T_{\mathrm{non\mbox{-}blank}},T_{\mathrm{edge}}$ & Visibility and edge thresholds fixed by Eq.~\eqref{eq:integer-thresholds}. \\
$O_i,L_i^r,\mathcal L_r$ & Append-only observation log of replica $i$, its prefix-consistent extension, and the $n-f$ extensions admitted in fragment $r$. \\
$\mathsf{Seq}_r[i],\mathsf{Head}_r[i]$ & Last committed log position and hash-chain head of replica $i$. \\
$\mathsf{TxRef}_r[u]$ & Canonical content commitment retained while $u$ is live; the corresponding bytes remain data-available. \\
$\mathsf{Pos}_r(i,u)$ & Unique committed first position of live transaction $u$ in replica $i$'s observation log. \\
$c_r(u),S_r(u)$ & Derived visibility count and state $\Blank$, $\Shaded$, or $\Solid$ of live transaction $u$. \\
$\mathsf{Live}_r,A_r$ & Admitted unfinalized set and the active subset induced by cumulative evidence. \\
$W_r(u,v)$ & Derived number of distinct replicas whose committed positions place $u$ before $v$. \\
$u\triangleleft v$ & Public total order over canonical transaction identifiers used only to break an exact pair-weight tie. \\
$\mathsf{EdgePred}_r,\widehat{\mathsf{EdgePred}}_r$ & Deterministic edge predicates for committed and candidate evidence states. \\
$\Sigma_r,\widehat\Sigma_r,\Sigma_r^{\mathsf{cand}}$ & Committed state, candidate pre-finalization state, and candidate post-finalization state. \\
$\mathsf{View}_r,\widehat{\mathsf{View}}_r$ & Deterministic derived views containing $c$, $W$, $S$, and $A$. \\
$h_r,\widehat h_r,D_r$ & Committed post-state identifier, candidate evidence-state identifier, and BFT-committed fragment digest. \\
$\mathsf{AuthCtx}$ & BFT-maintained authorized order-leader epoch and public key. \\
$G_r^\star,\widehat G_r^\star$ & Protocol-defined graphs of the committed and candidate pre-finalization states. \\
$G_{\mathrm{utig}},G_{\mathrm{utig},r},\widehat G_{\mathrm{utig},r}$ & Persistent UTIG and its installed and candidate graph materializations. \\
$\mathsf{TouchedNodes}_r,\mathsf{TouchedPairs}_r$ & Nodes and unordered pairs whose derived state or edge predicate may change in round $r$. \\
$\mathsf{Solid}_r,\mathsf{Anc}_r(u)$ & Solid active transactions and the ancestors of $u$ in the candidate active dependency graph. \\
$K_r^{\mathsf{safe}}$ & Maximal Solid subset that is down-closed in the candidate active dependency graph. \\
$F_r,\mathsf{Part}_r,b(u),\beta_r$ & Serialized finalized prefix, SCC/fairness-batch partition, assigned batch index, and largest committed batch index. \\
$\mathcal P_r,\mathcal F_r,\sigma_r^{\mathsf L}$ & Structural certificate, proof-carrying fragment, and order-leader signature. \\
$\mathsf{Done}_r,\mathsf{DoneRoot}_r,\mathsf{DoneRef}_r$ & Finalized-identifier dictionary, its commitment, and its canonical checkpoint reference. \\
$K_{\mathrm{cp}}$ & Number of committed fragments between checkpoints. \\
$\mathsf{Snap}_r,\mathsf{CP}_r,\mathsf{AvailQC}_r$ & Canonical live-evidence snapshot, committed checkpoint, and snapshot-availability certificate. \\
\hline
\end{tabular}
\end{table*}

%-------------------------------------------------------------------------------
\section{Expanded Protocol Procedures}
\label{app:expanded-protocol}
%-------------------------------------------------------------------------------

This appendix is the normative field-level specification of the procedures
summarized in Section~\ref{sec:protocol}.

\subsection{LocalOrder Admissibility}

An empty extension has $\ell_i^+=\ell_i^-$ and
$\eta_i^+=\eta_i^-$.  The transport limit $lo\text{-}size$ bounds a contiguous
chunk but never permits skipped positions.  A correct replica signs all
extension fields preceding $\sigma_i$, retransmits its next uncommitted chunk,
and does not emit an empty extension while a position is pending.  Conflicting
signed extensions of one committed position/head constitute equivocation
evidence.  Following a committed handoff that preserves the fragment sequence,
the replica re-signs the same next chunk with the successor epoch; its committed
sequence number and head do not change.

After GST, a correct order leader selects $n-f$ admissible senders using the
public rotating priority $\mathsf{prio}_r(i)=(i-r)\bmod n$.  Already finalized
identifiers remain part of hash-chain validation but require neither contents
nor new live evidence.  Every new non-finalized identifier must resolve to
available, application-valid contents with a unique content commitment.

\subsection{Evidence Application}

\begin{algorithm}[h!]
\caption{\textsc{ApplyEvidence}: Expanded Evidence Transition}
\label{alg:apply-evidence}
\footnotesize
\Input{$\Sigma_{r-1}$, $\mathsf{View}_{r-1}$, verified
       $(\mathsf{epoch}_r,r)$, and $\mathcal L_r$}
\Output{$\widehat\Sigma_r$, $\widehat{\mathsf{View}}_r$,
        $\mathsf{TouchedNodes}_r$, and $\mathsf{TouchedPairs}_r$, or reject}
\If{$\mathcal L_r$ lacks exactly $n-f$ distinct admissible senders}
   {\KwRet reject}
$\widehat h_r\gets\Commit(\mathsf{epoch}_r,r,\Sigma_{r-1}.h,
H(\Enc(\mathcal L_r));\mathsf{AUTIG/prestate})$\;
$\widehat\Sigma_r\gets\Sigma_{r-1}$;
set its version, epoch, and $h$ fields to
$r$, $\mathsf{epoch}_r$, and $\widehat h_r$\;
$\widehat{\mathsf{View}}_r\gets\mathsf{View}_{r-1}$;
$\mathsf{TouchedNodes}_r,\mathsf{TouchedPairs}_r\gets\varnothing$\;
\ForEach{$L_i^r\in\mathcal L_r$}{
  Verify its signature, epoch/sequence, predecessor position/head,
  contiguity, and uniqueness; for each new non-finalized identifier, verify
  content availability and application validity\;
  $K_i^-\gets\{u:\mathsf{Pos}_{r-1}(i,u)\ne\bot\}$;
  $N_i\gets\varnothing$\;
  \ForEach{$u$ at position $p$ in $\mathsf{chunk}_i^r$}{
    \eIf{$u\in\mathsf{Done}_{r-1}$}{
      Treat $u$ as continuity-only\;
    }{
      \If{$\mathsf{Pos}_{r-1}(i,u)\ne\bot$}{\KwRet reject}
      Install $\widehat{\mathsf{Pos}}_r(i,u)=p$ and store
      $\widehat{\mathsf{TxRef}}_r[u]$, or require equality with its retained
      commitment; add $u$ to
      $\widehat{\mathsf{Live}}_r$, $N_i$, and
      $\mathsf{TouchedNodes}_r$; increment $\widehat c_r(u)$\;
    }
  }
  $K_i^+\gets K_i^-\cup N_i$\;
  \ForEach{$\{u,v\}\subseteq K_i^+$ with
           $\{u,v\}\not\subseteq K_i^-$}{
    Increment exactly one of $\widehat W_r(u,v)$ and $\widehat W_r(v,u)$
    according to the unique positions; add $\{u,v\}$ to
    $\mathsf{TouchedPairs}_r$\;
  }
  $\widehat{\mathsf{Seq}}_r[i]\gets\ell_i^+$;
  $\widehat{\mathsf{Head}}_r[i]\gets\eta_i^+$\;
}
\ForEach{$u\in\mathsf{TouchedNodes}_r$}{
  Recompute $\widehat S_r(u)$ from $\widehat c_r(u)$\;
  \If{$u$ changes from $\Blank$ to non-Blank}{
    Add $u$ to $\widehat A_r$ and every live pair incident to $u$ to
    $\mathsf{TouchedPairs}_r$\;
  }
}
\KwRet $(\widehat\Sigma_r,\widehat{\mathsf{View}}_r,
         \mathsf{TouchedNodes}_r,\mathsf{TouchedPairs}_r)$
\end{algorithm}

\subsection{Safe Extraction}

\begin{algorithm}[h!]
\caption{\textsc{SafeExtract}: Expanded Safe Extraction}
\label{alg:safe-extract}
\footnotesize
\Input{$\widehat\Sigma_r$, $\widehat{\mathsf{View}}_r$, and
       $\widehat G_{\mathrm{utig},r}$}
\Output{$F_r$, $\mathsf{Part}_r$, and $\mathcal P_r$}
Compute the SCCs and condensation DAG of
$\widehat G_{\mathrm{utig},r}=\widehat G_r^\star$\;
Mark a component unsafe if it contains a non-Solid vertex or has an unsafe
predecessor; let $K_r^{\mathsf{safe}}$ be the union of the remaining components\;
Set $F_r$ to the safe components in minimum-$\mathsf{cid}$ Kahn order, with
identifiers increasing within each component\;
Construct $\mathsf{Part}_r$, inward/outward trees, and ranks for the included
components\;
Construct a predecessor map rooted at Shaded vertices outside
$K_r^{\mathsf{safe}}$: each non-root $v$ stores $p(v)$ with
$\widehat{\mathsf{EdgePred}}_r(p(v)\to v)$, covering every excluded Solid vertex\;
\KwRet $(F_r,\mathsf{Part}_r,\mathcal P_r)$
\end{algorithm}

\subsection{Fragment Verification}

\begin{algorithm}[h!]
\caption{\textsc{VerifyFragment}: Expanded Verification}
\label{alg:verify-fragment-expanded}
\footnotesize
\Input{Committed $\Sigma_{r-1}$; committed predecessor digest $D_{r-1}$;
       current $\mathsf{AuthCtx}$; candidate $\mathcal F_r$}
\Output{$(accept,\Sigma_r^{\mathsf{cand}},h_r)$ or reject}
Require the header's epoch and leader key to equal $\mathsf{AuthCtx}$; verify
the order-leader signature and Eq.~\eqref{eq:fragment}; require
$r=\Sigma_{r-1}.r+1$\;
Require $h_{r-1}=\Sigma_{r-1}.h$ and
$H_{\mathcal F,\prev}=D_{r-1}$\;
Derive $\mathsf{View}_{r-1}$ and run
\textsc{ApplyEvidence}; reject on failure or a mismatching $\widehat h_r$\;
\If{$F_r$ contains duplicates, finalized identifiers, or identifiers outside
    $\widehat A_r$}{\KwRet reject}
\If{$\mathrm{set}(F_r)\not\subseteq\mathsf{Solid}_r$}{\KwRet reject}
Require exact partition/tree/rank key coverage and reconstruct every edge of
$\widehat G_r^\star[\mathrm{set}(F_r)]$ from local $\widehat W_r$\;
For each claimed component, require every tree edge to be reconstructed, its
out-tree to reach every vertex from the root, and its in-tree to reach the root
from every vertex\;
Require every reconstructed cross-component edge to increase rank strictly;
thus the claimed components are the maximal SCCs of the induced graph\;
Recompute the public Kahn order, require $\mathsf{Rank}_r$ to encode that order,
and require $F_r$ to list SCCs by rank and identifiers increasingly within each
SCC\;
\ForEach{$(x,y)\in(\widehat A_r\setminus\mathrm{set}(F_r))
                   \times\mathrm{set}(F_r)$}{
  \If{$\widehat{\mathsf{EdgePred}}_r(x\to y)$ under local $\widehat W_r$}
     {\KwRet reject}
}
Require $\mathsf{BlockForest}_r$ to be a vertex-keyed map without
duplicates and with keys in
$\widehat A_r\setminus\mathrm{set}(F_r)$\;
Require each root to be Shaded, parentless, and at depth zero; for every
non-root $v$, require one parent $p(v)$,
$\widehat{\mathsf{EdgePred}}_r(p(v)\to v)$, and depth one greater than its
parent\;
Require parent-pointer traversal from every excluded Solid vertex to terminate
at a Shaded root and reject records outside those certified paths\;
$\Sigma_r^{\mathsf{cand}}\gets
  \textsc{Finalize}(\widehat\Sigma_r,F_r,\mathsf{Part}_r)$;
require $\Sigma_r^{\mathsf{cand}}.h=h_r$\;
\KwRet $(accept,\Sigma_r^{\mathsf{cand}},h_r)$
\end{algorithm}

%-------------------------------------------------------------------------------
\section{Running Example}
\label{app:running-example}
%-------------------------------------------------------------------------------

Consider $n=5$, $f=1$, and $\gamma=1$, which gives $q_h=4$,
$T_{\mathrm{edge}}=2$, and $T_{\mathrm{solid}}=3$.  Suppose the committed
prefixes expose
\[
 L_1=[e,a,b,c,d],\qquad
 L_2=[b,c,a,e,d],\qquad
 L_3=[c,a,b,d],
\]
while the remaining replicas have not committed positions for these
transactions.  Transactions $a,b,c,d$ are Solid and $e$ is Shaded.  The
derived graph contains the cycle $a\to b\to c\to a$, the edges
$a,b,c\to d$, and $e\to d$.  Hence the SCC $\{a,b,c\}$ is safe, whereas the
Shaded ancestor $e$ prevents $d$ from entering $K_r^{\mathsf{safe}}$.

The fragment contains the fairness batch $\{a,b,c\}$, serialized as
$[a,b,c]$.  Its certificate supplies two spanning trees for that SCC, its
condensation rank, and the blocker edge $e\to d$.  A follower derives the
corresponding edges from its authenticated positions and rejects if any active
vertex outside $\{a,b,c\}$ has an edge into the proposed set.  Once later
committed prefixes make $e$ Solid, the blocker disappears and $d$ becomes
eligible for a subsequent safe closure.

%-------------------------------------------------------------------------------
\section{Detailed Security Proofs}
\label{app:autig-security-details}
%-------------------------------------------------------------------------------

The proofs below use the candidate pre-finalization state
$\widehat\Sigma_r$; hats are omitted where the round is clear.

\begin{lemma}[Position uniqueness and deterministic weights]
\label{lem:position-uniqueness}
For any live transaction $u$, replica $i$ has at most one
$\mathsf{Pos}_r(i,u)$.  For every unordered live pair $\{u,v\}$, the two
weights are uniquely determined by Eq.~\eqref{eq:position-weight}, and one
replica contributes to at most one direction.
\end{lemma}

\begin{proof}
An admitted extension starts at the committed sequence number and log head,
and positions inside the extension are contiguous.  Algorithm~\ref{alg:apply-evidence}
rejects a second live position for the same sender and identifier.  The two
indicators for $u\prec v$ and $v\prec u$ in Eq.~\eqref{eq:position-weight}
are mutually exclusive, so the committed positions uniquely determine both
weights.
\end{proof}

\begin{proof}[Proof of Theorem~\ref{thm:state-agreement}]
We induct on the BFT-committed fragment sequence.  The claim holds for the
common genesis state.  Assume that two correct replicas have installed the
same $\Sigma_{r-1}$.  Agreement and data availability at the hosting BFT give
them the same committed fragment $\mathcal F_r$, including
$\mathcal L_r$, $F_r$, and $\mathsf{Part}_r$; its header fixes the same round
and order-leader epoch.  For every admitted LocalOrder,
signature, epoch, sequence, predecessor-head, continuity, uniqueness, and
content-validity checks are deterministic.  Thus
Algorithm~\ref{alg:apply-evidence} advances the same
$r$ and $\mathsf{epoch}_r$ state fields, advances the same
$\mathsf{Seq}_r$ and $\mathsf{Head}_r$, assigns the same unique
$\mathsf{Pos}_r$ entries, and obtains identical $\mathsf{Live}_r$ and
$\mathsf{TxRef}_r$.  The definition of $\mathsf{View}_r$ then gives the same
visibility counts, pair weights, transaction states, and active set.

Given the same accepted $F_r$ and $\mathsf{Part}_r$, \textsc{Finalize}
removes the same live entries, records the same round and fairness-batch
indices in $\mathsf{Done}_r$, and advances $\beta_r$ identically.  Canonical
encoding fixes $\mathsf{DoneRoot}_r$, while
the transition-identifier definitions fix $\widehat h_r$ and $h_r$.  Hence both
replicas install the same $\Sigma_r$ and derive the same
$\mathsf{View}_r$.  The induction proves the claim for every committed round;
uncommitted candidate states do not enter the induction.
\end{proof}

\begin{lemma}[Authenticated observation coverage]
\label{lem:committed-observation-coverage}
Let $H_{xy}$ contain at least $q_h$ correct replicas that received $x$ before
$y$, and suppose $x$ is not finalized in candidate pre-finalization state
$\widehat\Sigma_r$.  Define
\[
 \widehat Y_r(y)=\{i:\widehat{\mathsf{Pos}}_r(i,y)\ne\bot\},
 \qquad \widehat c_r(y)=|\widehat Y_r(y)|.
\]
Then
\begin{equation}
 \widehat W_r(x,y)\ge \widehat L_r(x,y)
   =\max\{0,q_h+\widehat c_r(y)-n\}.
 \label{eq:coverage-lower-bound}
\end{equation}
\end{lemma}

\begin{proof}
The two sets are drawn from $n$ replicas, so
$|H_{xy}\cap\widehat Y_r(y)|\ge q_h+\widehat c_r(y)-n$.  Every replica in the
intersection is correct and has an authenticated prefix containing $y$.
Prefix completeness and
the assumption that $x$ remains unfinalized imply that the same prefix contains
$x$ at a smaller position.  Each intersection member contributes one unit to
$\widehat W_r(x,y)$.
\end{proof}

\begin{lemma}[Fair edge formation]
\label{lem:fair-edge-formation}
Under Eq.~\eqref{eq:integer-feasibility}, if the premises of
Lemma~\ref{lem:committed-observation-coverage} hold and
$y\in\mathsf{Solid}_r$, then
$\widehat W_r(x,y)\ge T_{\mathrm{edge}}$ and
$\widehat W_r(x,y)>\widehat W_r(y,x)$.  Hence
$\widehat{\mathsf{EdgePred}}_r(x\to y)$ holds independently
of deterministic tie-breaking.
\end{lemma}

\begin{proof}
Write $c=\widehat c_r(y)$ and $L=q_h+c-n$.  Solidity gives
$c\ge n-2f$, and therefore
\[
 L\ge q_h-2f\ge n-q_h+1=T_{\mathrm{edge}},
\]
where the second inequality is Eq.~\eqref{eq:integer-feasibility}.  Every
contributor to $\widehat W_r(y,x)$ belongs to $\widehat Y_r(y)$, whereas the $L$ contributors
identified above cannot contribute the reverse direction.  Thus
$\widehat W_r(y,x)\le c-L$.  Finally,
\[
 2L-c\ge2q_h+c-2n\ge2q_h-n-2f\ge1,
\]
which proves
$\widehat W_r(x,y)\ge L>c-L\ge\widehat W_r(y,x)$.  Moreover, $y$ is Solid and
therefore belongs to $\widehat A_r$; since
$\widehat W_r(x,y)\ge T_{\mathrm{edge}}=T_{\mathrm{non\mbox{-}blank}}$,
$\widehat c_r(x)\ge T_{\mathrm{non\mbox{-}blank}}$ and
$x\in\widehat A_r$.  All conjuncts of
$\widehat{\mathsf{EdgePred}}_r(x\to y)$ therefore
hold.
\end{proof}

\begin{lemma}[Safe-closure certificate]
\label{lem:safe-closure-certificate}
If Algorithm~\ref{alg:verify-fragment-expanded} accepts the structural certificate and
frontier scan, then $\mathrm{set}(F_r)=K_r^{\mathsf{safe}}$ and the claimed
groups are exactly the SCCs of
$\widehat G_r^\star[\mathrm{set}(F_r)]$.
\end{lemma}

\begin{proof}
The inward and outward trees make every claimed component strongly connected.
Complete internal cross-component inspection and strictly increasing ranks
rule out a cycle among claimed components, establishing SCC maximality.  Every
included vertex is Solid.  The frontier scan rules out an edge from outside
$F_r$ into $F_r$; any outside ancestor path would contain such a first entering
edge.  Hence every ancestor of an included vertex is also included and Solid,
which proves $F_r\subseteq K_r^{\mathsf{safe}}$.  Conversely, the
unique-parent and depth checks give each excluded Solid vertex a finite
verified path from a Shaded root and therefore a non-Solid ancestor.  No
excluded Solid vertex belongs to
$K_r^{\mathsf{safe}}$, proving the reverse inclusion.  Recomputing the public
tie-breaking rules proves the partition and linearization claims.  Before
these structural checks, Algorithm~\ref{alg:apply-evidence} verifies content
binding, data availability, and application validity for every newly admitted
live identifier.  Previously admitted live identifiers remain bound by
$\mathsf{TxRef}_r$ to the contents validated when they were introduced.
Therefore every transaction in an accepted $F_r$ is valid and data-available,
which proves the remaining claim of Theorem~\ref{thm:safe-fragment}.
\end{proof}

\begin{proof}[Proof of Theorem~\ref{thm:batch-fairness}]
Consider the committed fragment that first finalizes $y$.  If $x$ was
finalized earlier, then $b(x)<b(y)$.  Otherwise,
Lemma~\ref{lem:fair-edge-formation} gives $x\to y$ in the candidate live
graph.  Theorem~\ref{thm:safe-fragment} makes $F_r$ down-closed, so
$x\in F_r$.  If $x$ and $y$ share an SCC, they receive the same fairness-batch
index; otherwise the condensation order places $x$ in an earlier batch.
Thus $b(x)\le b(y)$.
\end{proof}

\begin{proof}[Proof of Theorem~\ref{thm:pruning_safety}]
The authoritative committed positions uniquely determine visibility counts,
transaction states, the active set, and both directional weights for every
live pair.  Applying $\mathsf{EdgePred}_r$ to these values therefore reconstructs
exactly the edge set of the unpruned UTIG.  SCC decomposition, condensation,
safe-set marking, and deterministic linearization are pure functions of this
edge set.  Consequently, cache eviction changes neither the rematerialized
graph nor $K_r^{\mathsf{safe}}$.
Algorithm~\ref{alg:verify-fragment-expanded} checks the
resulting certificate against the same authoritative follower state, so both
executions produce the same acceptable fragment.  The argument does not apply
to deletion of an authoritative position belonging to a live transaction,
which soft pruning forbids.
\end{proof}

\begin{proof}[Proof of Theorem~\ref{thm:hard-deletion}]
Let $u$ be finalized in round $r$.  For any transaction $x$ whose
correct-replica receive-order observations satisfy the fairness antecedent
relative to $u$, the proof of
Theorem~\ref{thm:batch-fairness} shows that $x$ is finalized in the same or
an earlier fairness batch.  Therefore a later chunk cannot reveal an
unfinalized fairness predecessor of $u$.  Removing $u$'s live positions and
evicting its incident cached weights does not change any pair weight whose two
endpoints remain live, and $\mathsf{Done}_r$ prevents a delayed occurrence of $u$ from
reactivating it.

The accepted set $F_r$ is down-closed in the candidate active dependency graph.
Consequently, no vertex that remains active after finalization can reach a
vertex in $F_r$: otherwise down-closure would require that remaining vertex
to belong to $F_r$.  A finalized vertex therefore cannot lie on a directed
path between two remaining active vertices.  Deleting $u$ and its incident
edges preserves both reachability and SCC structure among all remaining
active transactions.  Together with the unchanged live--live pair values,
this shows that future safe closures and fairness-batch assignments are
unchanged.
\end{proof}

\begin{proof}[Proof of Theorem~\ref{thm:recovery}]
The checkpoint binds $\mathsf{AuthCtx}_r$, the canonical live-data references,
positions, log heads, $\mathsf{DoneRef}_r$, and its authenticated root.  Resolving the
reference reconstructs the finalized dictionary.  After the referenced data
are retrieved and validated, these values uniquely derive support counts, states,
the active set, and pair weights.  Starting from this common base, a committed
handoff deterministically updates only $\mathsf{AuthCtx}$, whereas a committed
fragment has one deterministic evidence transition and one deterministic
finalization transition and advances $D$.  Induction over the BFT-committed
AUTIG suffix yields the same $\Sigma_r$, $h_r$, $D_r$, and authorization context
as uninterrupted execution.
\end{proof}

%-------------------------------------------------------------------------------
\section{Incremental UTIG Materialization}
\label{app:incremental-materialization}
%-------------------------------------------------------------------------------

The protocol-defined graph for candidate state $\widehat\Sigma_r$ is
$\widehat G_r^\star=(\widehat A_r,\widehat E_r^\star)$, where
$(u,v)\in\widehat E_r^\star$ exactly when
$\widehat{\mathsf{EdgePred}}_r(u\to v)$ evaluates to true using the positions and
weights derived from $\widehat\Sigma_r$.  The materialized graph
$\widehat G_{\mathrm{utig},r}$ is a leader-side candidate cache; it is not an
additional source of protocol truth.

\begin{lemma}[Touched-pair completeness]
\label{lem:touched-pair-completeness}
For every unordered pair $\{u,v\}\subseteq\widehat A_r$, if its membership or
orientation in the protocol-defined graph can differ after
Algorithm~\ref{alg:apply-evidence}, then
$\{u,v\}\in\mathsf{TouchedPairs}_r$.
\end{lemma}

\begin{proof}
There are two cases.  First, if both endpoints were active before the
transition, an edge predicate can change only if a new committed position
changes $W(u,v)$ or $W(v,u)$.  The pair-update loop adds every pair receiving
such a first contribution, including changes that preserve the same threshold
side but alter relative support.  Second, suppose at least one endpoint was
not previously active.  The endpoint can enter $\widehat A_r$ only by receiving
a new committed position that moves it from $\Blank$ to a non-Blank state.
The state-update loop then adds every live incident pair.  If neither case
applies, the active membership and both directional weights are unchanged, so
the deterministic edge predicate is unchanged.
\end{proof}

\begin{proof}[Proof of Theorem~\ref{thm:incremental-equivalence}]
Assume inductively that the persistent UTIG before round $r$ equals the graph
derived from $\Sigma_{r-1}$.  Algorithm~\ref{alg:apply-evidence} does not modify
an authoritative position except by assigning a previously absent first
position.  By Lemma~\ref{lem:touched-pair-completeness}, every edge whose
existence or orientation can change is re-evaluated, while every untouched
edge has the same endpoints and predicate inputs as before.  Adding newly
active vertices and applying the refreshed predicates therefore yields
$\widehat G_{\mathrm{utig},r}=\widehat G_r^\star$ before extraction.  Once the
fragment commits,
finalization deletes exactly the vertices in $F_r$ and their incident edges.
It does not change a position or pair value whose endpoints both remain live.
The installed post-finalization materialization consequently satisfies
$G_{\mathrm{utig},r}=G_r^\star$, completing the induction.
\end{proof}

\begin{lemma}[Version-consistent extraction]
\label{lem:version-consistent-extraction}
If graph refresh, safe extraction, and certificate construction read one
immutable evidence version identified by $\widehat h_r$, then their output is
equivalent to a serial execution of Algorithms~\ref{alg:apply-evidence} and
\ref{alg:safe-extract} on $\widehat\Sigma_r$.
\end{lemma}

\begin{proof}
The immutable version fixes the active set, every pair value, and therefore
the complete protocol-defined graph.  By
Theorem~\ref{thm:incremental-equivalence}, the refreshed UTIG represents that
same graph.  SCC decomposition, safe-set marking, and deterministic
linearization are pure functions of this version.  Equation~\eqref{eq:fragment}
binds $\widehat h_r$ to the certificate and proposed prefix, and
Algorithm~\ref{alg:verify-fragment-expanded} reconstructs the same evidence version
before checking them.  Thus, an update to a later evidence version cannot be
observed partially by an accepted fragment.  An implementation may realize
this requirement with immutable snapshots, copy-on-write versions, or an
equivalent single-writer/multiple-reader protocol.
\end{proof}

%-------------------------------------------------------------------------------
\section{Pruning and State Retention}
\label{app:pruning-retention}
%-------------------------------------------------------------------------------

AUTIG distinguishes authoritative evidence from reconstructible graph state.
The authoritative state consists of the round and epoch, fairness-batch
counter, per-replica log heads and sequence numbers, live set, live-data
references, live first positions, finalized-identifier dictionary and its
root, and committed transition identifier, exactly as specified in
Eq.~\eqref{eq:follower-state}.
Visibility counts, transaction
states, aggregate weights, active-set indexes, dependency edges, SCCs, and
condensation metadata are deterministic derivatives.  This distinction gives
the following retention rules.

\begin{enumerate}
  \item Derived views and graph structures---including aggregate weights,
        active-set indexes, edges, adjacency indexes, SCCs, condensation
        ranks, and extraction workspaces---may be evicted and reconstructed
        from the authenticated evidence state.
  \item The live-data reference and authoritative positions of every live
        transaction remain available, independently of age, current state, or
        distance from the extracted prefix.
  \item After a BFT-committed fragment finalizes $u$, its live-data reference
        and live positions may be deleted, and every derived cache entry
        incident to $u$ may be evicted.  The committed log heads and sequence
        numbers remain, and $u$ remains queryable through $\mathsf{Done}$.
  \item No authoritative position belonging to an unfinalized transaction is
        hard-deleted; consequently, every required live--live pair value
        remains deterministically reconstructible.
\end{enumerate}

\begin{lemma}[Delayed-extension compatibility]
\label{lem:delayed-extension-compatibility}
Deleting live positions incident to a finalized transaction does not prevent
the admission of a later prefix-consistent LocalOrder that contains that
identifier.
\end{lemma}

\begin{proof}
Admissibility checks the extension's predecessor sequence number and hash head,
neither of which is deleted by finalization.  An occurrence found in
$\mathsf{Done}$ is consumed for log continuity but is not assigned a new live
position and contributes no incident weight.  Every non-finalized identifier
later in the same contiguous chunk is processed at its absolute log position.
Consequently, delayed delivery advances the replica log without reactivating
the finalized transaction or changing a live--live comparison.
\end{proof}

Theorem~\ref{thm:pruning_safety} follows because soft pruning removes only
deterministic derivatives.  Theorem~\ref{thm:hard-deletion} and
Lemma~\ref{lem:delayed-extension-compatibility} establish the stronger boundary
for committed finalized identifiers.  In contrast, deleting a live--live
position can change a cumulative weight, an edge predicate, and a future safe
closure; such a deletion is not a semantics-preserving pruning operation.
The state size is therefore bounded by the currently retained live evidence
and finalized-membership representation, rather than by a constant independent
of workload.

%-------------------------------------------------------------------------------
\section{Checkpoint Recovery and Order-Leader Handoff}
\label{app:checkpoint-handoff}
%-------------------------------------------------------------------------------

Recovery begins from the latest checkpoint whose identifier is committed by
the hosting BFT.  The snapshot reconstructs the complete authoritative
evidence state at that checkpoint; the committed AUTIG suffix supplies every
later handoff, evidence, and finalization transition.  The tagged reference
$\mathsf{DoneRef}_j$ canonically encodes either
$\mathsf{Inline}(\mathsf{Done}_j,\mathsf{DoneRoot}_j)$ or
$\mathsf{AppCheckpoint}(\mathsf{appcp}_j,\mathsf{DoneRoot}_j)$; the latter
retrieves the dictionary from the identified application checkpoint.

\begin{algorithm}[h!]
\caption{\textsc{RecoverState}: Checkpoint Recovery and Committed Replay}
\label{alg:recover-state}
\footnotesize
\Input{Committed checkpoint $\mathsf{CP}_j$; candidate snapshot
       $\mathsf{Snap}_j$; BFT-committed AUTIG suffix $\mathcal C_{>j}$}
\Output{Recovered committed context $(\Sigma,D,\mathsf{AuthCtx})$, or reject}
Verify that $\mathsf{CP}_j$ is BFT-committed and contains $n-f$ valid
availability signatures on $(\mathsf{AuthCtx}_j,j,h_j,D_j,d_j)$\;
Require canonical snapshot encoding,
$\Commit(\mathsf{Snap}_j;\mathsf{AUTIG/checkpoint\mbox{-}snapshot})=d_j$,
and equality of its round and transition identifier with $(j,h_j)$\;
Decode $\mathsf{Seq}_j$, $\mathsf{Head}_j$, $\mathsf{Live}_j$,
$\mathsf{TxRef}_j$, $\mathsf{Pos}_j$, $\mathsf{DoneRef}_j$, and $\beta_j$;
resolve $\mathsf{DoneRef}_j$ to
$(\mathsf{Done}_j,\mathsf{DoneRoot}_j)$\;
Verify $\mathsf{Live}_j\cap\mathsf{Done}_j=\varnothing$,
$\mathrm{dom}(\mathsf{TxRef}_j)=\mathsf{Live}_j$,
$\mathsf{DoneRoot}_j=\Commit(\mathsf{Done}_j;
\mathsf{AUTIG/finalized\mbox{-}set})$, and that every stored position is
consistent with its sender's committed sequence number\;
Resolve every $\mathsf{TxRef}_j[u]$, require its transaction digest to equal
$u$, and verify the deterministic application-validity predicate\;
$\Sigma\gets\Sigma_j$ reconstructed from the verified snapshot\;
$\mathsf{View}\gets\mathsf{Derive}(\Sigma)$; validate and discard any supplied
derived cache\;
$D\gets D_j$; $\mathsf{AuthCtx}\gets\mathsf{AuthCtx}_j$\;
\ForEach{record $R$ in $\mathcal C_{>j}$ in BFT-committed order}{
  \eIf{$R$ is a committed handoff to $\mathsf{AuthCtx}'$}{
    Require $R$ to extend
    $(\mathsf{AuthCtx},\Sigma.r+1,\Sigma.h,D)$ and name the next epoch;
    $\mathsf{AuthCtx}\gets\mathsf{AuthCtx}'$\;
  }{
    Parse $R$ as fragment $\mathcal F_k$ and require its committed digest and
    available bytes and referenced contents to match\;
    $(accept,\Sigma^{\mathsf{cand}},h_k)\gets
       \textsc{VerifyFragment}(\Sigma,D,\mathsf{AuthCtx},\mathcal F_k)$\;
    \If{verification rejects}{\KwRet reject}
    $\Sigma\gets\Sigma^{\mathsf{cand}}$; $D\gets H_{\mathcal F_k}$\;
  }
}
\KwRet $(\Sigma,D,\mathsf{AuthCtx})$
\end{algorithm}

\paragraph{Snapshot availability.}
Every $n-f$-signature $\mathsf{AvailQC}_j$ contains at least $n-2f$ correct
signers.  Each correct signer attests to the complete tuple
$(\mathsf{AuthCtx}_j,j,h_j,D_j,d_j)$ and retains the matching snapshot and referenced
recovery data until a successor checkpoint becomes committed and available.  Under
eventual communication,
at least one such holder supplies a digest-matching snapshot.  Collision
resistance prevents a different snapshot from satisfying the committed
digest.

\paragraph{Exact reconstruction.}
The checkpoint fixes the authorization context and the snapshot fixes the
authoritative base state.  Equation~\eqref{eq:position-weight}
then fixes every cumulative pair value, including external--external pairs
that no finalized prefix would reveal.  Each suffix handoff updates only the
authorized epoch and key; each suffix fragment contains the LocalOrders,
finalization data, and referenced contents required by
Algorithm~\ref{alg:verify-fragment-expanded}.  Replay excludes uncommitted
proposals and local timeout observations.
The induction in Theorem~\ref{thm:recovery} therefore yields the exact current
state, rather than a lower bound or conservative superset of historical
evidence.

\begin{corollary}[Order-leader handoff continuity]
\label{cor:handoff-continuity}
An order leader that starts from the committed context returned by
Algorithm~\ref{alg:recover-state} materializes the same current UTIG and
extracts the same safe closure as an uninterrupted correct order leader.
\end{corollary}

\begin{proof}
Theorem~\ref{thm:recovery} gives equality of the authoritative follower state,
predecessor digest, and authorization context, while
Theorem~\ref{thm:incremental-equivalence} makes the materialized UTIG a
deterministic derivative of that state.  The successor ignores uncommitted
candidate states.  A fragment from a replaced leader fails the recovered
authorization-context check and cannot alter the committed history.
\end{proof}
